\section{Stolz 公式}
\label{sec:1.6}

下面介绍一种求极限时非常重要而且使用方便的方法，它就是 Stolz 公式：
\[
\lim_{n\to+\infty} \frac{y_n}{x_n} = \lim_{n\to+\infty} \frac{y_n - y_{n-1}}{x_n - x_{n-1}}.
\]
当 $y_n = \sum\limits_{k=1}^n a_k$ 时，特别有效.

\begin{theorem}[$\frac{*}{\infty}$ 型 Stolz 公式]\label{thm:1.6.1}
设有数列 $\{x_n\}$，$\{y_n\}$，其中 $\{x_n\}$ 严格增，且 $\lim\limits_{n\to+\infty} x_n = +\infty$（注意：不必 $\lim\limits_{n\to+\infty} y_n = +\infty$）。如果
\[
\lim_{n\to+\infty} \frac{y_n - y_{n-1}}{x_n - x_{n-1}} = a\quad (\text{实数}, +\infty, -\infty),
\]
则
\[
\lim_{n\to+\infty} \frac{y_n}{x_n} = a = \lim_{n\to+\infty} \frac{y_n - y_{n-1}}{x_n - x_{n-1}}.
\]
\end{theorem}

\begin{proof}
(1) $a$ 为实数。

$\forall \varepsilon > 0$，因为 $\lim\limits_{n\to+\infty} \frac{y_n - y_{n-1}}{x_n - x_{n-1}} = a$，所以 $\exists N_1 \in \mathbb{N}$，当 $n > N_1$ 时，有
\[
\left|\frac{y_n - y_{n-1}}{x_n - x_{n-1}} - a\right| < \frac{\varepsilon}{2},
\]
即
\[
a - \frac{\varepsilon}{2} < \frac{y_n - y_{n-1}}{x_n - x_{n-1}} < a + \frac{\varepsilon}{2},
\]
\[
\left(a - \frac{\varepsilon}{2}\right)(x_n - x_{n-1}) < y_n - y_{n-1} < \left(a + \frac{\varepsilon}{2}\right)(x_n - x_{n-1}).
\]
类推有
\[
\left(a - \frac{\varepsilon}{2}\right)(x_{n-1} - x_{n-2}) < y_{n-1} - y_{n-2} < \left(a + \frac{\varepsilon}{2}\right)(x_{n-1} - x_{n-2}),
\]
\[
\cdots\cdots
\]
\[
\left(a - \frac{\varepsilon}{2}\right)(x_{N_1+1} - x_{N_1}) < y_{N_1+1} - y_{N_1} < \left(a + \frac{\varepsilon}{2}\right)(x_{N_1+1} - x_{N_1}).
\]
上面各式相加得到
\[
\left(a - \frac{\varepsilon}{2}\right)(x_n - x_{N_1}) < y_n - y_{N_1} < \left(a + \frac{\varepsilon}{2}\right)(x_n - x_{N_1}),
\]
\[
a - \frac{\varepsilon}{2} < \frac{y_n - y_{N_1}}{x_n - x_{N_1}} < a + \frac{\varepsilon}{2}.
\]
对\textbf{固定}的 $N_1$，因为 $\lim\limits_{n\to+\infty} x_n = +\infty$，所以，$\exists N > N_1$，s.t. 当 $n > N$ 时，有
\[
\left|\frac{y_{N_1} - a x_{N_1}}{x_n}\right| < \frac{\varepsilon}{2},\quad 0 < \frac{x_{N_1}}{x_n} < 1.
\]
于是，
\begin{align*}
\left|\frac{y_n}{x_n} - a\right| &= \left|\frac{y_n - y_{N_1}}{x_n} + \frac{y_{N_1} - a x_{N_1}}{x_n} - a\left(1 - \frac{x_{N_1}}{x_n}\right)\right|\\
&= \left|\frac{y_{N_1} - a x_{N_1}}{x_n} + \left(1 - \frac{x_{N_1}}{x_n}\right)\left(\frac{y_n - y_{N_1}}{x_n - x_{N_1}} - a\right)\right|\\
&\leqslant \left|\frac{y_{N_1} - a x_{N_1}}{x_n}\right| + \left|\frac{y_n - y_{N_1}}{x_n - x_{N_1}} - a\right| < \frac{\varepsilon}{2} + \frac{\varepsilon}{2} = \varepsilon.
\end{align*}
这就证明了
\[
\lim_{n\to+\infty} \frac{y_n}{x_n} = a.
\]

(2) $a = +\infty$。

因为 $\lim\limits_{n\to+\infty} \frac{y_n - y_{n-1}}{x_n - x_{n-1}} = a = +\infty$，所以 $\exists N \in \mathbb{N}$，当 $n > N$ 时，$\frac{y_n - y_{n-1}}{x_n - x_{n-1}} > 1$，$y_n - y_{n-1} > x_n - x_{n-1} > 0$，即 $\{y_n\}$ 严格增。又由于
\begin{align*}
y_n - y_N &= (y_n - y_{n-1}) + (y_{n-1} - y_{n-2}) + \cdots + (y_{N+1} - y_N)\\
&> (x_n - x_{n-1}) + (x_{n-1} - x_{n-2}) + \cdots + (x_{N+1} - x_N)\\
&= x_n - x_N,
\end{align*}
根据 $\lim\limits_{n\to+\infty} x_n = +\infty$，知 $\lim\limits_{n\to+\infty} y_n = +\infty$。应用 (1) 的结果得到
\[
\lim_{n\to+\infty} \frac{x_n}{y_n} = \lim_{n\to+\infty} \frac{x_n - x_{n-1}}{y_n - y_{n-1}} = \lim_{n\to+\infty} 1 \bigg/ \frac{y_n - y_{n-1}}{x_n - x_{n-1}} = 0.
\]
于是，
\[
\lim_{n\to+\infty} \frac{y_n}{x_n} = \lim_{n\to+\infty} 1 \bigg/ \frac{x_n}{y_n} = +\infty = \lim_{n\to+\infty} \frac{y_n - y_{n-1}}{x_n - x_{n-1}}.
\]

(3) $a = -\infty$。

由 (2) 知，
\begin{align*}
\lim_{n\to+\infty} \frac{-y_n}{x_n} &= \lim_{n\to+\infty} \frac{(-y_n) - (-y_{n-1})}{x_n - x_{n-1}}\\
&= -\lim_{n\to+\infty} \frac{y_n - y_{n-1}}{x_n - x_{n-1}} = +\infty,
\end{align*}
即
\[
\lim_{n\to+\infty} \frac{y_n}{x_n} = -\lim_{n\to+\infty} \frac{-y_n}{x_n} = -\infty = \lim_{n\to+\infty} \frac{y_n - y_{n-1}}{x_n - x_{n-1}}.
\]
\end{proof}

\begin{remark}\label{rem:1.6.1}
当 $\lim\limits_{n\to+\infty} \frac{y_n - y_{n-1}}{x_n - x_{n-1}} = \infty$ 时，$\{x_n\}$ 严格增且 $\lim\limits_{n\to+\infty} x_n = +\infty$ 时，并不能推出 $\lim\limits_{n\to+\infty} \frac{y_n}{x_n} = \infty$。

反例：$x_n = n, y_n = \left[1 + (-1)^n\right]n^2$，此时，
\[
\lim_{n\to+\infty} \frac{y_n - y_{n-1}}{x_n - x_{n-1}} = \infty,
\]
但 $\lim\limits_{n\to+\infty} \frac{y_n}{x_n} \neq \infty$。
\end{remark}

\begin{theorem}[$\frac{0}{0}$ 型 Stolz 公式]\label{thm:1.6.2}
设 $\{x_n\}$ 严格减，且 $\lim\limits_{n\to+\infty} x_n = 0$，$\lim\limits_{n\to+\infty} y_n = 0$。若
\[
\lim_{n\to+\infty} \frac{y_n - y_{n-1}}{x_n - x_{n-1}} = a\quad (\text{实数}, +\infty, -\infty),
\]
则
\[
\lim_{n\to+\infty} \frac{y_n}{x_n} = a = \lim_{n\to+\infty} \frac{y_n - y_{n-1}}{x_n - x_{n-1}}.
\]
\end{theorem}

\begin{proof}
(1) $a$ 为实数。

$\forall \varepsilon > 0$，因为 $\lim\limits_{n\to+\infty} \frac{y_n - y_{n+1}}{x_n - x_{n+1}} = \lim\limits_{n\to+\infty} \frac{y_n - y_{n-1}}{x_n - x_{n-1}} = a$，所以 $\exists N \in \mathbb{N}$，当 $n > N$ 时，有
\[
a - \frac{\varepsilon}{2} < \frac{y_n - y_{n+1}}{x_n - x_{n+1}} < a + \frac{\varepsilon}{2},\quad x_n - x_{n+1} > 0,
\]
\[
\left(a - \frac{\varepsilon}{2}\right)(x_n - x_{n+1}) < y_n - y_{n+1} < \left(a + \frac{\varepsilon}{2}\right)(x_n - x_{n+1}).
\]
类似定理~\ref{thm:1.6.1}(1) 的证明有
\[
\left(a - \frac{\varepsilon}{2}\right)(x_n - x_{n+p}) < y_n - y_{n+p} < \left(a + \frac{\varepsilon}{2}\right)(x_n - x_{n+p}).
\]
令 $p \to +\infty$，则由 $x_{n+p} \to 0$，$y_{n+p} \to 0$，得到
\[
\left(a - \frac{\varepsilon}{2}\right)x_n \leqslant y_n \leqslant \left(a + \frac{\varepsilon}{2}\right)x_n.
\]
由于 $x_n > 0$，有
\[
a - \varepsilon < a - \frac{\varepsilon}{2} \leqslant \frac{y_n}{x_n} \leqslant a + \frac{\varepsilon}{2} < a + \varepsilon,
\]
所以，
\[
\lim_{n\to+\infty} \frac{y_n}{x_n} = a.
\]

(2) $a = +\infty$。

$\forall A > 0$，因为 $\lim\limits_{n\to+\infty} \frac{y_n - y_{n+1}}{x_n - x_{n+1}} = \lim\limits_{n\to+\infty} \frac{y_n - y_{n-1}}{x_n - x_{n-1}} = +\infty$，所以 $\exists N \in \mathbb{N}$，当 $n \geqslant N$ 时，有
\[
\frac{y_n - y_{n+1}}{x_n - x_{n+1}} > 2A,
\]
类似上述论证有
\[
y_n - y_{n+p} > 2A(x_n - x_{n+p}).
\]
令 $p \to +\infty$，由 $y_{n+p} \to 0$，$x_{n+p} \to 0$，得到
\[
y_n \geqslant 2A x_n,\quad \frac{y_n}{x_n} \geqslant 2A > A,
\]
所以，
\[
\lim_{n\to+\infty} \frac{y_n}{x_n} = +\infty.
\]

(3) $a = -\infty$。

类似 (2) 的证明或将 (2) 的结论应用到 $\{-y_n\}$ 即得。
\end{proof}

现在我们给出应用 Stolz 公式的典型例题。

\begin{example}\label{ex:1.6.1}
设 $\lim\limits_{n\to+\infty} a_n = a$（实数，$+\infty, -\infty$），则 $\lim\limits_{n\to+\infty} \frac{a_1 + a_2 + \cdots + a_n}{n} = a$。
\end{example}

\begin{proof}
\textit{证法 1}\quad 应用 $\varepsilon\text{-}N, A\text{-}N$ 法，参阅例~\ref{ex:1.1.15}。

\textit{证法 2}\quad 令 $y_n = a_1 + a_2 + \cdots + a_n$，$x_n = n$，则由 Stolz 公式得到
\begin{align*}
\lim_{n\to+\infty} \frac{a_1 + a_2 + \cdots + a_n}{n} &= \lim_{n\to+\infty} \frac{(a_1 + a_2 + \cdots + a_n) - (a_1 + a_2 + \cdots + a_{n-1})}{n - (n - 1)}\\
&= \lim_{n\to+\infty} \frac{a_n}{1} = \lim_{n\to+\infty} a_n = a.
\end{align*}
\end{proof}

\begin{example}\label{ex:1.6.2}
设 $\lim\limits_{n\to+\infty} a_n = a$，则 $\lim\limits_{n\to+\infty} \frac{a_1 + 2a_2 + \cdots + n a_n}{n^2} = \frac{a}{2}$（$a$ 为实数，$+\infty, -\infty$）。
\end{example}

\begin{proof}
\textit{证法 1}\quad 应用 $\varepsilon\text{-}N, A\text{-}N$ 法，参阅练习题~1.1 之第~9 题。

\textit{证法 2}\quad 因为 $\lim\limits_{n\to\infty} a_n = a$，根据极限定义易证数列：
\[
a_1, a_2, a_2, a_3, a_3, a_3, \cdots, \underbrace{a_n, \cdots, a_n}_{n \text{ 个}}, \cdots
\]
的极限也为 $a$。应用例~\ref{ex:1.1.15}，有
\begin{align*}
\lim_{n\to+\infty} \frac{a_1 + 2a_2 + \cdots + n a_n}{n^2} &= \lim_{n\to+\infty} \frac{a_1 + 2a_2 + \cdots + n a_n}{\frac{n(n+1)}{2}} \cdot \frac{\frac{n(n+1)}{2}}{n^2}\\
&= \lim_{n\to+\infty} \frac{a_1 + 2a_2 + \cdots + n a_n}{\frac{n(n+1)}{2}} \cdot \frac{1}{2}\left(1 + \frac{1}{n}\right)\\
&= a \cdot \frac{1}{2}(1 + 0) = \frac{a}{2}.
\end{align*}

\textit{证法 3}\quad 
\begin{align*}
\lim_{n\to+\infty} \frac{a_1 + 2a_2 + \cdots + n a_n}{n^2} &\xlongequal{\text{Stolz 公式}} \lim_{n\to+\infty} \frac{n a_n}{2n - 1}\\
&= \lim_{n\to+\infty} \frac{1}{2 - \frac{1}{n}} a_n = \frac{a}{2}.
\end{align*}
\end{proof}

\begin{example}\label{ex:1.6.3}
设 $k$ 为自然数。证明：
\begin{enumerate}
    \item 
    \[
    \lim_{n\to+\infty} \frac{1^k + 2^k + \cdots + n^k}{n^{k+1}} = \frac{1}{k+1};
\]
    \item 
    \[
    \lim_{n\to+\infty} n\left(\frac{1^k + 2^k + \cdots + n^k}{n^{k+1}} - \frac{1}{k+1}\right) = \frac{1}{2}.
\]
\end{enumerate}
\end{example}

\begin{proof}
(1) 
\begin{align*}
\lim_{n\to+\infty} \frac{1^k + 2^k + \cdots + n^k}{n^{k+1}} &\xlongequal{\text{Stolz 公式}} \lim_{n\to+\infty} \frac{n^k}{n^{k+1} - (n-1)^{k+1}}\\
&\xlongequal{\text{二项式定理}} \lim_{n\to+\infty} \frac{n^k}{C_{k+1}^1 n^k - C_{k+1}^2 n^{k-1} + \cdots + (-1)^{k+1}}\\
&= \frac{1}{C_{k+1}^1} = \frac{1}{k+1}.
\end{align*}

(2) 
\begin{align*}
&\lim_{n\to+\infty} n\left(\frac{1^k + 2^k + \cdots + n^k}{n^{k+1}} - \frac{1}{k+1}\right)\\
=&\; \lim_{n\to+\infty} \frac{(k+1)(1^k + 2^k + \cdots + n^k) - n^{k+1}}{(k+1)n^k}\\
\xlongequal{\text{Stolz 公式}}&\; \lim_{n\to+\infty} \frac{(k+1)n^k - \left[n^{k+1} - (n-1)^{k+1}\right]}{(k+1)\left[n^k - (n-1)^k\right]}\\
\xlongequal{\text{二项式定理}}&\; \lim_{n\to+\infty} \frac{\frac{k(k+1)}{2} n^{k-1} + \cdots}{(k+1)\cdot k n^{k-1} + \cdots} = \frac{1}{2}.
\end{align*}
\end{proof}

\begin{remark}\label{rem:1.6.2}
当 $k \in (-1, +\infty)$ 但 $k$ 不为自然数时，不能应用二项式定理。但例~\ref{ex:1.6.3} 的结论仍成立，这就要用到后面的知识：
\begin{enumerate}
    \item 可由函数极限 $\lim\limits_{x\to 0} \frac{(1+x)^{k+1} - 1}{x} = k+1$ 推得，或用定积分定义及 Newton-Leibniz 公式得到。
    \item 可由 $(1+x)^\alpha\;(\alpha > 0)$ 在 $x = 0$ 处的 Taylor 展开求得。
\end{enumerate}
\end{remark}

\begin{example}\label{ex:1.6.4}
设级数 $\sum\limits_{n=1}^\infty a_n$ 收敛，又 $\{p_n\}$ 为严格增的正数数列，且 $\lim\limits_{n\to+\infty} p_n = +\infty$。证明：
\[
\lim_{n\to+\infty} \frac{p_1 a_1 + \cdots + p_n a_n}{p_n} = 0.
\]
\end{example}

\begin{proof}
令 $S_n = a_1 + \cdots + a_n, n \in \mathbb{N}$，由 $\sum\limits_{n=1}^\infty a_n$ 收敛，记 $\lim\limits_{n\to+\infty} S_n = S$。于是，$a_1 = S_1, a_n = S_n - S_{n-1}, n = 2, 3, \cdots$，且
\begin{align*}
\lim_{n\to+\infty} \frac{p_1 a_1 + \cdots + p_n a_n}{p_n} &= \lim_{n\to+\infty} \frac{p_1 S_1 + p_2(S_2 - S_1) + \cdots + p_n(S_n - S_{n-1})}{p_n}\\
&= \lim_{n\to+\infty} \left[\frac{S_1(p_1 - p_2) + S_2(p_2 - p_3) + \cdots + S_{n-1}(p_{n-1} - p_n)}{p_n} + S_n\right]\\
&\xlongequal{\text{Stolz 公式}} \lim_{n\to+\infty} \left[\frac{S_n(p_n - p_{n+1})}{p_{n+1} - p_n} + S_n\right]\\
&= \lim_{n\to+\infty} (-S_n + S_n) = \lim_{n\to+\infty} 0 = 0.
\end{align*}
\end{proof}

\begin{exblock}{练习题 1.6}
    \item 设 $C_n^k = \frac{n!}{k!(n-k)!}$ 为组合数。应用 Stolz 公式证明：
    \[
    \lim_{n\to+\infty} \frac{\sum\limits_{k=0}^n \ln C_n^k}{n^2} = \frac{1}{2}.
\]

    \item 应用 Stolz 公式证明：
    \begin{enumerate}
        \item $\lim\limits_{n\to+\infty} \frac{\sum\limits_{k=1}^n \sqrt{k}}{n^{\frac{3}{2}}} = \frac{2}{3}$;
        \item $\lim\limits_{n\to+\infty} n \left[\frac{\sum\limits_{k=1}^n \sqrt{k}}{n^{\frac{3}{2}}} - \frac{2}{3}\right] = \frac{1}{2}$.
    \end{enumerate}
\end{exblock}

\begin{exblock}{思考题 1.6}
    \item 设 $0 < x_1 < 1$，$x_{n+1} = x_n(1 - x_n)$，$n = 1, 2, \cdots$。证明：$\lim\limits_{n\to+\infty} n x_n = 1$。进而设 $0 < x_1 < \frac{1}{q}$，其中 $0 < q \leqslant 1$，并且 $x_{n+1} = x_n(1 - q x_n)$，$n \in \mathbb{N}$。证明：$\lim\limits_{n\to+\infty} n x_n = \frac{1}{q}$。

    \item 由 Toeplitz 定理导出 $\frac{\infty}{\infty}$ 型的 Stolz 公式。

    \item 设数列 $\{a_n\}$ 满足 $\lim\limits_{n\to+\infty} a_n \sum\limits_{i=1}^n a_i^2 = 1$。证明：$\lim\limits_{n\to+\infty} \sqrt[3]{3n} a_n = 1$。
\end{exblock}
